\chapter{Conclusion and Future Work}
\label{ch:conclusion}

In this report we studied stack allocation of objects in Java, an optimization that reduces both
the access cost of objects and the load on the garbage collector. We surveyed two staged
static\,+\,JIT approaches that make the JIT compiler's escape analysis more precise without
making it slower: PYE, whose conditional values let a static analysis summarize methods despite
missing library code, and optimistic stack allocation with dynamic heapification, which acts on
static escape results at run time and repairs the cases where dynamic features such as class
loading invalidate them. The latter reports a 71\% increase in stack allocations, a 5.3\%
reduction in GC time and an 8.8\% reduction in execution time.

We then used OpenJ9's debug counters to look at how stack allocation behaves on the DaCapo
benchmarks. The 3000\,B per-method budget for stack-allocated objects is far from exhausted: no
compiled method used more than 6.4\% of it, so the size-based eviction heuristic never came into
play. Arrays, on the other hand, rarely reach the stack. Only \StackRateOfCand\% of candidate
arrays are stack allocated, and \TopReasonShare\% of the arrays rejected as candidates are
rejected because their size is not a compile-time constant.

\section{Future Work}

These observations suggest the following research directions.

\paragraph{Arrays of dynamic size (\TopReasonShare\% of rejections).}
\begin{itemize}
  \item Use profiling information to find arrays whose size turns out to be constant, or bounded,
        at run time, and reserve stack-frame space for them. A method is typically recompiled
        several times as it gets hotter, so profile data collected during earlier executions is
        available to later, more optimized compilations.
  \item Add heapification checks wherever such an array may be reallocated, so that the program
        stays correct when a run-time size does not fit the reserved space.
  \item Statically identify arrays that are never reallocated, and skip the heapification checks
        for them.
\end{itemize}

\paragraph{Large arrays.} Split large arrays into chunks in the style of arraylets (a spine that
points to fixed-size leaves), so that the chunks fit into the stack-frame budget. The trade-off is
one extra indirection on every element access.

\paragraph{The candidate-to-stack gap.} Most candidate arrays still end up on the heap. Finding
out where in escape analysis they are lost, and why, may well be the larger opportunity, and it
is a prerequisite for knowing how much the two directions above can gain.
